%% toc: P2 21 21.2,21.3
\section{Craig interpolation with free variables}
\label{sec:p02-w02-p02-037-craig}

\emph{Purpose.}  Theorem~\ref{thm:p02-craig-interpolation} in
\cref{sec:p02-craig-beth} states interpolation for sentences.  Applications
with parameters, such as definitions of \(n\)-ary predicates and
interpolation between formulas met inside a derivation, need the formula
form, in which the interpolant must also keep its free variables among the
variables the two formulas share.  This section states that form with both
of its conditions, gives the procedure that reduces it to the sentence case,
and names the complete proofs of the sentence case elsewhere in the book.

\emph{Inputs and conclusion.}  The input is a pair of first-order formulas
\(\varphi,\psi\), over compatible signatures, such that
\(\varphi\to\psi\) is valid.  The conclusion is a formula \(\chi\) between
them whose nonlogical symbols are common to \(\varphi\) and \(\psi\) and
whose free variables are free in both.  Signatures, terms, formulas and
satisfaction are those of \cref{sec:fol-syntax-semantics-01},
\cref{sec:fol-syntax-semantics-02} and \cref{sec:fol-semantics-02}:
equality and \(\bottom\) are logical, structures have nonempty domains, and
a formula is valid when it is satisfied in every structure under every
assignment.

\begin{hbconvention}{def:p02-w02-p02-037-vocabulary}{Vocabulary and compatible signatures}
For a formula \(\varphi\), the \emph{occurrence vocabulary}
\(\operatorname{voc}(\varphi)\) is the set of function symbols (constants
included) and relation symbols that occur in \(\varphi\), each with its kind
and arity.  The equality sign and \(\bottom\) are logical tokens, not
signature symbols, and never belong to \(\operatorname{voc}(\varphi)\).
Signatures \(\Lang_1\) and \(\Lang_2\) are \emph{compatible} when every
symbol belonging to both has the same kind and the same arity in both; then
their union \(\Lang_1\cup\Lang_2\) is again a signature.  For
\(\varphi\in\operatorname{Fm}(\Lang_1)\) and \(\psi\in\operatorname{Fm}(\Lang_2)\), the statement
\(\entails\varphi\to\psi\) means validity over all
\((\Lang_1\cup\Lang_2)\)-structures.  By \cref{lem:reduct-invariance} the
same statement holds over every signature extending \(\Lang_1\cup\Lang_2\),
and every structure for a smaller signature containing
\(\operatorname{voc}(\varphi)\cup\operatorname{voc}(\psi)\) expands to a
\((\Lang_1\cup\Lang_2)\)-structure (extra functions constant at one chosen
element, extra relations empty), so the choice of the ambient signature
does not matter.
\end{hbconvention}

\begin{hbtheorem}{thm:p02-w02-p02-037-craig-formula}{proved=annex:p02-w02-p02-037-craig-formula}{Craig interpolation, formula form}
Let \(\Lang_1\) and \(\Lang_2\) be compatible signatures, let
\(\varphi\in\operatorname{Fm}(\Lang_1)\) and \(\psi\in\operatorname{Fm}(\Lang_2)\), and suppose
\(\entails\varphi\to\psi\).  Then there is a formula
\(\chi\in\operatorname{Fm}(\Lang_1\cup\Lang_2)\) such that
\[
 \entails\varphi\to\chi,
 \qquad
 \entails\chi\to\psi,
\]
\[
 \operatorname{voc}(\chi)\subseteq
   \operatorname{voc}(\varphi)\cap\operatorname{voc}(\psi),
 \qquad
 \FV(\chi)\subseteq\FV(\varphi)\cap\FV(\psi).
\]
Equality and \(\bottom\) may occur in \(\chi\) whatever the two
vocabularies are.  If \(\varphi\) and \(\psi\) are sentences, \(\chi\) is a
sentence; this case is Theorem~\ref{thm:p02-craig-interpolation}.
\index{Craig interpolation theorem!formula form}
\end{hbtheorem}

\emph{Reading the two conditions.}  The vocabulary condition does not bound
the free variables: if \(\chi\) is an interpolant, so is
\(\chi\land z=z\) for any variable \(z\), and its vocabulary is unchanged
because equality is logical.  The free-variable condition is therefore a
separate requirement.  It says that an interpolant may speak about the
shared parameters of \(\varphi\) and \(\psi\) and about nothing else
outside its own quantifiers.  When the two formulas share no nonlogical
symbol and no free variable, \(\chi\) is a sentence built from equality
and \(\bottom\) alone; examples are \(\bottom\), \(\neg\bottom\), and
cardinality sentences such as \(\exists x\,\exists y\,\neg x=y\).

\begin{hbmethod}{def:p02-w02-p02-037-reduction}{From a formula implication to a sentence implication and back}
Given \(\entails\varphi\to\psi\) as in
Theorem~\ref{thm:p02-w02-p02-037-craig-formula}:
\begin{enumerate}
 \item List \(V=\FV(\varphi)\cup\FV(\psi)=\{v_1,\ldots,v_k\}\), with the
       \(v_i\) pairwise distinct.
 \item Choose pairwise distinct constant symbols \(c_1,\ldots,c_k\) that
       belong to neither \(\Lang_1\) nor \(\Lang_2\).
 \item Form the sentences
       \[
        \varphi^{\mathrm{c}}=\varphi[c_1/v_1]\cdots[c_k/v_k],
        \qquad
        \psi^{\mathrm{c}}=\psi[c_1/v_1]\cdots[c_k/v_k].
       \]
       A constant contains no
       variable, so these substitutions rename no bound variable and
       replace exactly the free occurrences of each \(v_i\).
 \item Take a sentence interpolant \(\theta\) for
       \(\entails\varphi^{\mathrm{c}}\to\psi^{\mathrm{c}}\) whose vocabulary lies in
       \(\operatorname{voc}(\varphi^{\mathrm{c}})\cap\operatorname{voc}(\psi^{\mathrm{c}})\).
 \item Rename the bound variables of \(\theta\) so that no quantifier binds
       a member of \(V\).
 \item Replace every occurrence of each \(c_i\) by \(v_i\); the result is
       \(\chi\).
\end{enumerate}
Side conditions.  The constants must be fresh and pairwise distinct as
symbols, but no axiom makes them denote distinct elements, so assignments
giving two free variables the same value are covered.  Step~5 must precede
step~6: for \(\theta=\forall v_1\,R(c_1,v_1)\), replacing \(c_1\) first
would give \(\forall v_1\,R(v_1,v_1)\), which captures the new occurrence;
renaming first gives \(\forall z\,R(c_1,z)\) and then
\(\forall z\,R(v_1,z)\).  The constant \(c_i\) occurs in
\(\varphi^{\mathrm{c}}\) exactly when \(v_i\in\FV(\varphi)\), and likewise for
\(\psi^{\mathrm{c}}\), so the vocabulary condition of step~4 is what yields the
free-variable condition of the theorem.
\end{hbmethod}

\emph{Proof routes.}  Annex~\ref{annex:p02-w02-p02-037-craig-formula}
carries out the procedure in full, with the substitution, coincidence and
reduct lemmas of Part~II, and takes the sentence case of step~4 from either
of two complete proofs in the book.
\begin{itemize}
 \item Theorem~\ref{thm:craig-interpolation}, proved in
       Annex~\ref{annex:p7.interp.3}: Maehara's partition method on
       cut-free derivations gives the relational sentence case, and a
       relational translation of functions, constants and equality extends
       it to arbitrary signatures.  That theorem also states the
       free-variable clause; the annex here writes out the passage from
       sentences to formulas step by step.
 \item Theorem~\ref{thm:p03-craig-robinson-beth}, proved in
       Annex~\ref{annex:p3.abstract.5}: sentence interpolation follows from
       the common-sentence form of Robinson joint consistency,
       Theorem~\ref{thm:p03-robinson-joint-consistency}, proved in
       Annex~\ref{annex:p3.abstract.1}, together with compactness.  This
       route uses compactness for arbitrary signatures, which the book states
       under the ambient Choice assumption (\cref{thm:fol-compactness}).
\end{itemize}

\begin{hbexample}{def:p02-w02-p02-037-craig-example}{An interpolant with one shared parameter}
Let \(R\) be binary and \(Q,S\) unary, and put
\[
 \varphi(x,y)=R(x,y)\land\forall z\,(R(x,z)\to Q(z)),
 \qquad
 \psi(y,w)=Q(y)\lor S(w).
\]
Any structure and assignment satisfying \(\varphi\) give the pair
\((s(x),s(y))\) to \(R\), and the universal conjunct, read at the value
\(s(y)\) for \(z\), gives \(Q\) at \(s(y)\); hence \(\psi\) holds and
\(\entails\varphi\to\psi\).  Here
\(\operatorname{voc}(\varphi)\cap\operatorname{voc}(\psi)=\{Q\}\) and
\(\FV(\varphi)\cap\FV(\psi)=\{y\}\).

Following Method~\ref{def:p02-w02-p02-037-reduction} with
\(V=\{x,y,w\}\) and fresh constants \(c_x,c_y,c_w\) gives
\[
 \varphi^{\mathrm{c}}=R(c_x,c_y)\land\forall z\,(R(c_x,z)\to Q(z)),
 \qquad
 \psi^{\mathrm{c}}=Q(c_y)\lor S(c_w).
\]
The common vocabulary of these sentences is \(\{Q,c_y\}\), and
\(\theta=Q(c_y)\) is a sentence interpolant: \(\varphi^{\mathrm{c}}\) forces
\(Q\) at the value of \(c_y\), which makes the first disjunct of
\(\psi^{\mathrm{c}}\) true.  No bound variable needs renaming, and replacing
\(c_y\) by \(y\) gives \(\chi=Q(y)\).  It satisfies both implications,
uses only \(Q\), and has \(\FV(\chi)=\{y\}\).

The atom \(R(x,y)\) also follows from \(\varphi\) but violates the
vocabulary condition, and \(Q(y)\land x=x\) satisfies both implications
and the vocabulary condition but violates the free-variable condition.
\end{hbexample}

\begin{hblemma}{thm:p02-w02-p02-037-equality-needed}{proved=annex:p02-w02-p02-037-equality-needed}{Equality is needed in interpolants}
Let \(P\) be unary, and put
\[
 \varphi=\exists x\,\exists y\,(P(x)\land\neg P(y)),
 \qquad
 \psi=\exists x\,\exists y\,\neg x=y.
\]
Then \(\entails\varphi\to\psi\) and
\(\operatorname{voc}(\varphi)\cap\operatorname{voc}(\psi)=\varnothing\),
but no formula \(\chi\) in which equality does not occur and with
\(\operatorname{voc}(\chi)=\varnothing\) satisfies both
\(\entails\varphi\to\chi\) and \(\entails\chi\to\psi\).
\end{hblemma}

\begin{hbwarning}{def:p02-w02-p02-037-craig-limits}{What the formula form does not give}
\begin{itemize}
 \item Lemma~\ref{thm:p02-w02-p02-037-equality-needed} shows that the
       exemption for equality is part of the theorem.  In a language
       without logical equality, such as the sequent calculus in which
       Maehara proves interpolation, the formula \(\psi\) of the lemma is
       not available.  For that setting Maehara's statement carries a
       separate clause: if the two formulas share no predicate symbol, then
       the negation of the first or the second itself is provable
       \cite[\S3, p.~236]{Maehara1960}.  This section proves only the
       form with logical equality.
 \item The theorem concerns validity over all structures with nonempty
       domains.  It states nothing about interpolation relative to a
       restricted class of structures.
 \item It states no bound on the length or quantifier rank of \(\chi\), no
       polarity condition on the occurrences of shared relation symbols
       (Lyndon's refinement), and no uniqueness: in the example above, both
       \(Q(y)\) and \(Q(y)\land\forall z\,(Q(z)\to Q(z))\) are interpolants.
\end{itemize}
\end{hbwarning}

\section{Beth definability: the renamed-copy formulation and its proofs}
\label{sec:p02-w02-p02-037-beth}

\emph{Purpose.}  Definition~\ref{def:p02-beth-definability} expresses
implicit definability as a consequence relation: the theory together with
its copy \(T(P')\) entails that \(P\) and \(P'\) agree.  The complete proofs
of Beth's theorem in the book use a model formulation instead: any two
models of the theory with the same reduct to the old signature interpret
the predicate identically (Theorem~\ref{thm:beth} and
Definition~\ref{def:p03-implicit-explicit-definability}).  This section
proves that the two formulations coincide, derives
Theorem~\ref{thm:p02-beth-definability} in its stated form from the proved
theorems, and records which proof each route uses.

\emph{Inputs and conclusion.}  The input is a signature \(\Lang\), an
\(n\)-ary relation symbol \(P\notin\Lang\) with \(n\geq0\), an
\(\Lang(P)\)-theory \(T(P)\) (a set of sentences,
Definition~\ref{def:fol-theory-satisfiability}) and pairwise distinct
variables \(x_1,\ldots,x_n\).  The signature and the theory may be
infinite, and \(\Lang\) may contain function symbols and constants.  The
conclusion is that the implicit and explicit conditions of
Definition~\ref{def:p02-beth-definability} are equivalent, with an
explicit definition already entailed by a finite part of \(T(P)\).  For
\(n=0\) the variable list and its quantifier block are empty.

\begin{hblemma}{thm:p02-w02-p02-037-renaming}{proved=annex:p02-w02-p02-037-renaming}{Renaming the predicate}
Let \(P'\notin\Lang(P)\) be a fresh \(n\)-ary relation symbol.  For an
\(\Lang(P)\)-formula \(\alpha\), write \(\alpha'\) for the result of
replacing every occurrence of \(P\) by \(P'\).  For an
\(\Lang(P,P')\)-structure \(\mathcal N\), let \(\mathcal N^{\mathrm{ren}}\) be the
\(\Lang(P)\)-structure with the domain and \(\Lang\)-interpretations of
\(\mathcal N\) and with \(P^{\mathcal N^{\mathrm{ren}}}=P'^{\mathcal N}\).  Then
for every \(\Lang(P)\)-formula \(\alpha\) and assignment \(s\),
\[
 \mathcal N,s\entails\alpha'
 \quad\Longleftrightarrow\quad
 \mathcal N^{\mathrm{ren}},s\entails\alpha .
\]
\end{hblemma}

\begin{hblemma}{thm:p02-w02-p02-037-implicit-agree}{proved=annex:p02-w02-p02-037-implicit-agree}{Implicit definability: two formulations agree}
In the setting above, the following are equivalent.
\begin{enumerate}
 \item \(T(P)\) implicitly defines \(P\) over \(\Lang\) in the sense of
       Definition~\ref{def:p02-beth-definability}:
       \(T(P)\cup T(P')\entails\forall x_1\cdots\forall x_n\,
       (P(x_1,\ldots,x_n)\leftrightarrow P'(x_1,\ldots,x_n))\).
 \item Any two models \(\mathcal M_1,\mathcal M_2\) of \(T(P)\) with
       \(\mathcal M_1\mathbin{\upharpoonright}\Lang=\mathcal M_2\mathbin{\upharpoonright}\Lang\) satisfy
       \(P^{\mathcal M_1}=P^{\mathcal M_2}\); this is the hypothesis of
       Theorem~\ref{thm:beth} and the implicit condition of
       Definition~\ref{def:p03-implicit-explicit-definability}.
\end{enumerate}
\end{hblemma}

\begin{hbcorollary}{thm:p02-w02-p02-037-beth}{proved=annex:p02-w02-p02-037-beth}{Beth definability, renamed-copy form}
Assume the Axiom of Choice in the ambient set theory, as
Theorem~\ref{thm:beth} and Theorem~\ref{thm:p03-beth-definability} do.
With \(\Lang\), \(P\), \(T(P)\) and \(x_1,\ldots,x_n\) as above,
\(T(P)\) implicitly defines \(P\) over \(\Lang\) if and only if it
explicitly defines \(P\) over \(\Lang\), both in the sense of
Definition~\ref{def:p02-beth-definability}.  Moreover, if \(T(P)\)
implicitly defines \(P\), there are a finite \(T_0\subseteq T(P)\) and an
\(\Lang\)-formula \(\chi\) with \(\FV(\chi)\subseteq\{x_1,\ldots,x_n\}\)
such that
\[
 T_0\entails\forall x_1\cdots\forall x_n\,
 \bigl(P(x_1,\ldots,x_n)\leftrightarrow\chi\bigr).
\]
The direction from explicit to implicit definability uses no choice.
\end{hbcorollary}

\emph{Proof routes.}  The corollary is
Theorem~\ref{thm:p02-beth-definability} with its proof supplied:
Annex~\ref{annex:p02-w02-p02-037-beth} combines
Lemma~\ref{thm:p02-w02-p02-037-implicit-agree} with either of two complete
proofs of the model formulation.
\begin{itemize}
 \item Theorem~\ref{thm:beth}, proved in Annex~\ref{annex:p7.interp.4}:
       compactness reduces the hypothesis to one finite implication between
       sentences with fresh constants for the variables, and Craig
       interpolation (Theorem~\ref{thm:craig-interpolation}) supplies the
       definition.  This is the constant device of
       Method~\ref{def:p02-w02-p02-037-reduction} applied to the
       \(n\)-tuple of arguments.  The conclusion there is a derivation
       \(T\derives\cdots\); soundness (\cref{thm:fol-soundness}) turns it
       into the entailment used here.
 \item Theorem~\ref{thm:p03-beth-definability}, proved in
       Annex~\ref{annex:p3.abstract.2}: Robinson joint consistency
       (Theorem~\ref{thm:p03-robinson-joint-consistency},
       Annex~\ref{annex:p3.abstract.1}) applied to the theory and its
       renamed copy gives the definition directly as a semantic
       consequence.
\end{itemize}
A third route uses neither of these theorems.
Annex~\ref{annex:p02-w02-p02-037-beth} also derives the forward direction
directly from Theorem~\ref{thm:p02-w02-p02-037-craig-formula}.  Compactness
gives a finite \(T_0\subseteq T(P)\) whose conjunction \(\alpha\) satisfies
\[
 \entails\bigl(\alpha\land P(x_1,\ldots,x_n)\bigr)\to
 \bigl(\alpha'\to P'(x_1,\ldots,x_n)\bigr).
\]
The two sides share only symbols of \(\Lang\) and the variables
\(x_1,\ldots,x_n\), so an interpolant in the sense of
Theorem~\ref{thm:p02-w02-p02-037-craig-formula} is an \(\Lang\)-formula
with free variables among \(x_1,\ldots,x_n\).
Lemma~\ref{thm:p02-w02-p02-037-renaming} turns the second implication back
into a statement about \(P\), and the two implications together give the
explicit definition.  This is Maehara's derivation of Beth's theorem from
Craig's \cite[\S4, p.~236]{Maehara1960}.  Maehara works with provability in
Gentzen's sequent calculus: he uses the finiteness of a derivation where the
annex uses compactness, and he replaces the copy by the original predicate
throughout a derivation where the annex uses the renaming lemma.  His
statement of Beth's theorem is the renamed-copy formulation of
Definition~\ref{def:p02-beth-definability} for a one-place predicate, with
derivability in place of entailment \cite[\S2, p.~236]{Maehara1960}.

Theorem~\ref{thm:p03-craig-robinson-beth}, proved in
Annex~\ref{annex:p3.abstract.5}, connects the first two routes: under
compactness, Craig interpolation and the common-sentence form of joint
consistency imply each other, and each implies Beth definability.

\begin{hbexample}{def:p02-w02-p02-037-beth-example}{Running the renamed-copy test}
Let \(\Lang=\{Q\}\) with \(Q\) unary, and let
\[
 T(P)=\bigl\{\forall x\,(P(x)\lor Q(x)),\quad
             \forall x\,(P(x)\to\neg Q(x))\bigr\}.
\]
From \(T(P)\) we get \(\forall x\,(P(x)\leftrightarrow\neg Q(x))\): the
first sentence gives \(\neg Q\to P\) and the second gives
\(P\to\neg Q\).  The same reasoning in the copy \(T(P')\) gives
\(\forall x\,(P'(x)\leftrightarrow\neg Q(x))\), so
\(T(P)\cup T(P')\entails\forall x\,(P(x)\leftrightarrow P'(x))\), and
\(P\) is implicitly defined.  The explicit definition is
\(\chi=\neg Q(x_1)\), with \(\FV(\chi)=\{x_1\}\), and the finite set
\(T_0=T(P)\) already entails it.

Drop the first sentence and keep only
\(T_1(P)=\{\forall x\,(P(x)\to\neg Q(x))\}\).  The one-element structures
with \(Q\) empty and \(P\) either empty or the whole domain are models of
\(T_1(P)\) with the same \(\Lang\)-reduct and different interpretations of
\(P\).  By Lemma~\ref{thm:p02-w02-p02-037-implicit-agree}, \(T_1(P)\) does
not implicitly define \(P\), and by the easy direction of
Corollary~\ref{thm:p02-w02-p02-037-beth} it has no explicit definition
either.
\end{hbexample}

\begin{hbwarning}{def:p02-w02-p02-037-beth-limits}{Scope of the statement}
\begin{itemize}
 \item The statement concerns a relation symbol, as
       Definition~\ref{def:p02-beth-definability} does.  Definability of a
       new function symbol is not stated here.
 \item The explicit definition is an \(\Lang\)-formula with no parameters
       beyond \(x_1,\ldots,x_n\); it may use equality.  An explicit
       definition in a larger language does not by itself give an
       interpolant in the common language; the nullary example after
       Theorem~\ref{thm:p03-craig-robinson-beth} in
       \cref{sec:p03-beth-craig} shows why.
 \item The forward direction as proved here uses compactness for
       arbitrary signatures and hence the ambient Choice assumption; no
       weaker base theory is claimed.
\end{itemize}
\end{hbwarning}

\emph{Relationships and sources.}  The sentence forms in
\cref{sec:p02-craig-beth} and the lookup rows of
Table~\ref{tab:p02-fol-metatheorems} are the statements whose formula
form and proof routes are given here.  The interpolation theorem is due to
Craig \cite{Craig1957} and the definability theorem to Beth
\cite{Beth1953}; the joint consistency route goes back to Robinson
\cite{RobinsonJointConsistency1956}.  No locator in these three papers is
relied on.  Maehara's note \cite{Maehara1960} counts free object variables
among the parameters of a formula, so his statement of interpolation already
restricts the free variables of the interpolant to those the two formulas
share \cite[\S\S2--3, p.~236]{Maehara1960}; this is the free-variable
condition of Theorem~\ref{thm:p02-w02-p02-037-craig-formula}.  His proof
of interpolation partitions the antecedent and succedent of a sequent and
proceeds by induction on a cut-free derivation
\cite[\S5, pp.~236--237]{Maehara1960}, which is the method of
Annex~\ref{annex:p7.interp.3}.  The proofs referred to in this section are
the book's own.
